\section{Conclusions}
In the summary session, we proudly believe that we have solved a big problem for the Washington State Department of Agriculture. 

First of all, we use the cellular automata model to simulate the geographical distribution of the pests, and conclude that if the pest are not prevented and controlled in time, the consequences will be raging. At the same time, the accuracy of our simulation model can reach an area of $150m \times 150m$.

Secondly, we combine FCEM and AHP to calculate the possibility of other types of insects being misclassified as pests. The results show that the eastern cicada killer, European hornet and horn tail are the three most likely insect to be misclassified visually.

Thirdly, we use LR to filter the reports, and the filtered reports get a comprehensive score based on TOPSIS from the four dimensions. And finally we get the priority ranking of the sighting reports. In order to fully test the accuracy of this system, we randomly selected samples to test and get the desired results.

Fourthly, we establish a model update mechanism from two aspects: new pest invasion points and human intervention prevention and control. And we determine the update frequency based on existing data to be updated once every new intrusion point appears and once a year. 

Finally, we use the current number of existing pests in each cell as a measure and set three indicators for eradicating pests. And we test the sensitivity of the optimized model to prove that our model is stable.

\vspace{1cm}